\scanpage{34}\mspage{33}
% L1–3
\deleted{\unsure{7|8}}8) $\forall\,X\in\mathcal M$, let $X'$ be its ``complement'' \deleted{\unsure{8|6}} Since $\{X,X'\}$ tame, then $\{\overline X,\overline{X'}\}$ \deleted{$\mathrm{Fr}\,X$ is \unsure{finite} $\unsure{\{\overline X,\overline{X'}\}|\mathrm{Fr}\,X}\in\mathcal M$} \inserted{also, one can} form
\[
\left\{\begin{aligned}
\mathrm{Fr}\,X &\overset{\text{def}}{=} \overline X\wedge\overline{X'}\\
(\mathrm{Fr}\,X)' &= X^{\circ}\vee X^{\prime\circ}
\end{aligned}\right.
\]
% slanted
\margin{property special to the case considered of the model of \unsure{forms|closed sets} \uline{of} \uline{dimension $1$}}
% L4–5
That said $\{X,\deleted{\unsure{X'}}\unsure{Y|Y'}\}$ tame \deleted{\unsure{if and only if}} $\{\mathrm{Fr}\,X,\mathrm{Fr}\,Y\}$ is tame $\Longrightarrow$ $\mathrm{Fr}\,X\cup\mathrm{Fr}\,X$ is tame.

% L6–10
\deleted{\uline{NB}} \deleted{Let $X$ be \unsure{closed|tame}, then} \deleted{\inserted{One says that $X$ \unsure{($\in\mathcal M$)|(of $\mathcal M$)} is \uline{nowhere dense} if}} \deleted{if $X^\circ=\varnothing$ i.e. $\overline{X'}=1$ ($X'$ \uline{dense}). In the case at hand, $X$ is \unsure{nowhere dense|one} if and only if $X$ is finite (as a subset of $L$), or again if and only if $\mathrm{Fr}\,X$}

% L11–12
In the case at hand, equivalent conditions on $X\in\mathcal M$
% L13 (a)
\[
\begin{aligned}
\text{a)}\quad & X\subset\deleted{\unsure{\mathrm{Fr}(X)}}X^{\bullet}
&&\left(\begin{aligned}&\text{i.e.}\quad X^{\circ}=\varnothing\quad\&\ \text{i.e.}\quad\overline{X'}=1\\[-2pt]&\hspace{8em}\text{i.e. }X'\text{ ``dense''}\end{aligned}\right)\\
% L14 (b)
\text{b)}\quad & X=X^{\bullet}
&&\left(\begin{aligned}&\text{i.e. }X\text{ closed nowhere dense \quad i.e. }X'\text{ open}\\[-2pt]&\hspace{12em}\text{dense}\end{aligned}\right)
% L15 (c)
\end{aligned}
\]
c) $X$ is finite as a subset of $L$

\begin{center}\rule{1.6cm}{0.45pt}\end{center}

% L16–20
Consider an ordered set $\mathcal M$ satisfying conditions 1) to 7) (with 4) strengthened in the margin as has been said~; in algebraic notation, we have
% L21–24
\diagram{\begin{tikzcd}[ampersand replacement=\&,row sep=0.1ex,column sep=1.8em,
  cells={nodes={anchor=base west}},execute at end picture={
  \draw ([xshift=-0.9em]alJoin.north west) -- ([xshift=-0.9em]alOrder.south west);
  \draw[->] ([xshift=0.1em]alOrder.north west)
      .. controls ([xshift=-0.45em]alMeet.west) and ([xshift=-0.7em]alJoin.west)
      .. ([xshift=-0.45em]alJoin.north west)
      .. controls (alJoin.north west) and ([xshift=0.5em]alJoin.north west)
      .. ([xshift=1.1em,yshift=0.25em]alJoin.north west);
  \draw ([xshift=2.6em]alOrder.north west) to[bend right=8] (alOrder.north west);
  \draw (alOrder.south west) to[bend right=8] ([xshift=2.6em]alOrder.south west);
}]
X+Y=(X\vee Y)\wedge(X\wedge Y)' \& |[alias=alJoin]| X\vee Y=X+Y-XY\deleted{\unsure{Y|X}}\\
X\cdot Y=X\wedge Y \& |[alias=alMeet]| X\wedge Y=X\cdot Y\\
0=\text{least element} \& |[alias=alOrder]| X\leq Y\Longleftrightarrow X\cdot Y=X\\
1=\text{great\textsuperscript{est} element} \&
\end{tikzcd}}
